\begin{frame}{Промежуточные выводы}
	\begin{itemize}
		\item Erlang реализует отказоустойчивость через изоляцию процессов и супервизоры;
		\item Ошибки не предотвращаются полностью, а локализуются и приводят к перезапуску;
		\item Система продолжает работу даже при сбоях отдельных компонентов;
		\item Подход особенно эффективен в распределённых и real-time системах.
	\end{itemize}
	\vspace{0.8em}
	\textbf{Главная идея:} стабильность достигается за счёт контролируемых (управляемых) локальных сбоев и быстрого восстановления.
\end{frame}